\section{Saran}
Untuk memperoleh gambaran spasial yang lebih detail, penelitian selanjutnya disarankan untuk
\begin{enumerate}
    \item Menambah jumlah titik pengamatan,
    \item Menggunakan data curah hujan dari stasiun BMKG atau radar cuaca untuk meningkatkan akurasi,
    \item Mempertimbangkan menguji model Spatio Temporal modern berbasis Machine Learning seperti LSTM Spatio Temporal, Graph Neural Network,
    \item Mengembangkan sistem Prediksi berbasis aplikasi untuk sistem peringatan dini berbasis STARIMA,
    \item Pembuatan Dashboard monitoring curah hujan real time
\end{enumerate}
Saran - saran diatas diharapkan dapat menjadi acuan bagi peneliti selanjutnya dalam mengembangkan model prediksi curah hujan yang lebih akurat, responsive, dan aplikatif.
Kombinasi pendekatan statistik klasik dan metode modern berbasis kecerdasan buatan dapat menghasilkan model spatio temporal yang jauh lebih akurat dalam memetakan dinamika iklim di Kawasan Perkotaan seperti Surabaya.